\documentclass{techbrief}
\title{Schrödinger equation is equivalent to information conservation}
\author{Gabriele Carcassi}
\category{Quantum mechanics}
\tags{Schrödinger equation, entropy, unitary evolution}
\begin{document}
	\maketitle
	\begin{abstract}We show that the Schrödinger equation is equivalent to preservation of the von Neumann entropy. That is, the time evolution of a quantum system follows the Schrödinger equation if and only if it preserves information entropy.\end{abstract}
\section{Introduction}

In quantum mechanics, the time evolution for an isolated system is postulated to follow the Schrödinger equation (i.e. unitary evolution). Here we show that, once the basic structure of quantum theory is assumed, the same equation is equivalent to requiring the conservation of information entropy. Conceptually, this states that the time evolution is deterministic and reversible if and only if the amount of information given to identify the initial state is exactly the same as the information given to identify the final state, which is an information theoretic way to characterize determinism and reversibility.

\section{Conditions}

The following condition represents our base theory.
\begin{condition}[Quantum states]{QST:H}
	The state space of a quantum system is represented by the projective space $\mathcal{P}(\mathcal{H})$ of a separable complex Hilbert space $\mathcal{H}$. The ensemble space is represented by the density operators $\mathcal{D}(\mathcal{H})$. The conditional probability is given by the Born rule $p(\psi|\phi) = \frac{\< \phi | \psi \>\< \psi | \phi \>}{\< \psi | \psi \>\< \phi | \phi \>}$. The entropy is given by the von Neumann entropy: $S(\rho) = -\tr(\rho \log \rho)$.
\end{condition}

Within our base theory, we can derive the following.
\begin{coro}[Born rule is entropy]\label{QM-ENT-PROB}
	Let $\psi, \phi \in \mathcal{P}(\mathcal{H})$ be two quantum states. The probability of measuring $\phi$ having prepared $\psi$ is uniquely determined by the entropy of their equal mixture and vice versa. More precisely, $p(\phi|\psi) = \lambda$ if and only if $S\left(\frac{1}{2} \rho_{\psi} + \frac{1}{2} \rho_{\phi} \right) =  - \frac{1+\sqrt{\lambda}}{2} \log \frac{1+\sqrt{\lambda}}{2} - \frac{1-\sqrt{\lambda}}{2} \log \frac{1-\sqrt{\lambda}}{2}.$
\end{coro}
The expression of the entropy as a function of $\lambda$ is recovered with a straight calculation. The function is strictly decreasing and therefore invertible.

Time evolution in quantum mechanics is given by the following.
\begin{condition}[Schrödinger equation]{DR-SCEQ}
	The evolution follows the equation $\frac{d}{dt} \psi(t) = \frac{H(t)}{\imath \hbar} \psi(t)$, where $H(t)$ is a self-adjoint operator.
\end{condition}

We will want to show that the previous condition is equivalent to the following conditions.
\begin{condition}[Differentiable evolution]{EVDIFF}
	The evolution is differentiable.
\end{condition}

\begin{condition}[Conservation of information entropy]{DR-INFO}
	Future and past ensembles have the same information entropy. That is, if $\rho_I$ is an initial ensemble, $\rho_F$ is the corresponding evolved future ensemble and $\rho_P$ is the corresponding reconstructed past ensemble, then $S(\rho_I) = S(\rho_F) = S(\rho_P)$.
\end{condition}

Note that \ref{DR-SCEQ} respects the following condition:
\begin{condition}[Evolution preserves statistical mixtures]{[B]EVMIX}
	The evolution of a statistical mixture is the statistical mixture of the evolutions. That is, let $\mathcal{D}$ be the ensemble space of the system, the evolution from each $t_0$ to each $t_1$, with $t_1\geq t_0$, is given by an affine map $\Phi_{t_0 \to t_1}:\mathcal{D}\to \mathcal{D}$ such that $\Phi_{t_0 \to t_0}=I$, $\Phi_{t_1 \to t_2}\circ \Phi_{t_0 \to t_1}=\Phi_{t_0 \to t_2}$ and $(\rho,t_0,t_1)\mapsto \Phi_{t_0 \to t_1}(\rho)$ is jointly continuous.
\end{condition}
Physically, the process is such that the output of a statistical mixture is the statistical mixture of the outputs. Since this is a physical requirement, we will take it as part of our base theory.

\section{Theorem}
\begin{thrm}[Schrödinger equation is information preservation]
	Over \ref{QST:H}+\ref{[B]EVMIX}, conditions \ref{DR-SCEQ} and \ref{DR-INFO}+\ref{EVDIFF} are equivalent.
\end{thrm}
\begin{proof}
	\ref{DR-SCEQ} $\implies$ \ref{DR-INFO}+\ref{EVDIFF}. The Schrödinger equation gives rise to a differentiable unitary propagator $U_{t_0\to t_1}$. Therefore, the evolved and reconstructed ensembles are respectively $\rho_F=U_{t_0\to t_1}\rho_IU_{t_0\to t_1}^\dagger$ and $\rho_P=U_{t_0\to t_1}^\dagger\rho_FU_{t_0\to t_1}$. Since unitary conjugation preserves eigenvalues, and the von Neumann entropy depends only on eigenvalues, $S(\rho_I)=S(\rho_F)=S(\rho_P)$.
	
	For \ref{DR-INFO}+\ref{EVDIFF} $\implies$ \ref{DR-SCEQ}, consider the evolution $\Phi_{t_0\to t_1}:\mathcal{D}(\mathcal{H})\to\mathcal{D}(\mathcal{H})$. To show that it is induced by a unitary map on $\mathcal{H}$, we proceed in steps.
	
	\emph{Step 1: Pure states correspond to pure states.} Pure states are the unique ensembles with $S(\rho)=0$. Since entropy is preserved across evolved and reconstructed ensembles, pure initial states correspond bijectively to pure future states.
	
	\emph{Step 2: The map on vectors is unitary.} Let $\psi_F$ and $\phi_F$ be the future states corresponding to any two initial states $\psi_I$ and $\phi_I$. Consider their equal mixtures. By \ref{[B]EVMIX}, $\Phi_{t_0\to t_1}(\frac{1}{2}\rho_{\psi_I}+\frac{1}{2}\rho_{\phi_I})=\frac{1}{2}\rho_{\psi_F}+\frac{1}{2}\rho_{\phi_F}$. By \ref{DR-INFO}, the two mixtures have the same entropy. Since by \ref{QM-ENT-PROB} the entropy of an equal mixture of two pure states is an invertible function of their conditional probability, $p(\psi_F|\phi_F)=p(\psi_I|\phi_I)$ for all pairs of states. The bijective map on pure states therefore preserves the modulus of the inner product. By Wigner's theorem, it is induced by either a unitary or an anti-unitary map $U_{t_0\to t_1}$ on $\mathcal{H}$. Joint continuity and $\Phi_{t_0\to t_0}=I$ exclude the anti-unitary case, so $U_{t_0\to t_1}$ is unitary.
	
	\emph{Step 3: Unitary implies Schrödinger.} Since the evolution is differentiable, $U_{t_0\to t_0+dt}=I+A(t_0)dt+o(dt)$. Unitarity implies $A(t_0)=-A(t_0)^\dagger$. Setting $A(t_0)=\frac{H(t_0)}{i\hbar}$ with $H(t_0)=H(t_0)^\dagger$ recovers the Schrödinger equation.
\end{proof}

\section{Conclusion}

We have seen that the Schrödinger equation is equivalent to the preservation of von Neumann entropy. This means that the dynamical law of quantum mechanics can be understood as the requirement that time evolution neither creates nor destroys information. This parallels the classical case where Hamiltonian evolution is equivalent to the preservation of Shannon entropy (Liouville's theorem) and of the definitional independence of degrees of freedom.

\end{document}
